\begin{abstract}
Protecting information in power infrastructure requires release decisions that account for relations between operational data. Prices, forecasts and schedules support transparency, but their combinations may reveal quantities that remain restricted or have not yet been released. This paper develops an audit for assessing that background-dependent disclosure risk at field level. The audit measures the largest increase in inferability produced by adding a field to an admissible background, separates this increase into single-field risk and background amplification, and identifies the combinations responsible. An efficient estimator scans candidate backgrounds, while dedicated retraining verifies selected evidence. A controlled dispatch case shows how output identifies a unit's offer segment and how prices then reveal a private markup; congestion changes which prices are informative. Evaluation on RTS-GMLC, NEM, PJM and CAISO data covers ten targets and 124,425 exhaustively retrained field sets. The median average-to-maximum contextual gain ratio is 0.32, indicating substantial variation across backgrounds. Single-field screening recovers only 13--17\% of critical fields. The proposed scan and verification procedure recovers 76--91\% without false positives in the evaluated settings. In a withholding application, strategies based on the estimated combination structure leave only 0.3--1.6\% of minimal unsafe combinations intact, compared with 95.5--98.2\% for single-field screening. The resulting profiles provide interpretable evidence for reviewing co-release conditions, release timing and subsequent reclassification of power-system information.
\end{abstract}
